%% !TEX root = ../manual_gaolv5.tex
% Chapter 3.
%% Chapter 3: Using GAOL. Sources: gaol.tex of f5bf180 lines 840-1108,
%% 1750-1905, 1906-2300 and 5789-5880; doc/using.md lines 1-227, 291-450,
%% 572-630 and 695-1069; examples/tutorial.md (the names of the examples).
\chapter{Using GAOL}\label{chap:overview}
%%=========================

The previous chapter built and installed GAOL. This chapter shows how a
program uses it: what it includes, how it is compiled and linked with the
flags of interval arithmetic that the inline operations of GAOL require, and
which of the two namespaces it opens. It then reviews what GAOL gives to the
algorithms of interval computing, the errors most often made with intervals,
and the names of the other interval libraries next to those of GAOL.

\medskip
In this chapter, we assume that GAOL is installed, and that its library and
its headers are accessible to the compiler (Chapter~\ref{chap:installation}).
The first section walks through a complete program, from the header it
includes to the enclosure it prints. The second section lists the headers a
program may include, and the third gives the flags of interval arithmetic and
the way to take them from CMake, from pkg-config and from meson. The fourth
and the fifth sections present the namespaces \code{gaol} and
\code{gaol\_ieee1788}, and the power, where they differ. The sixth section
tells how an integer is taken, and the seventh what the compiler says of a
result thrown away. The eighth section gathers the tools of interval
algorithms, the ninth summarizes how GAOL reports an error, the tenth reviews
the common errors, and the last one gives the names of IBEX, Codac, C-XSC,
PROFIL/BIAS and Boost.Interval next to those of GAOL.

\section{A first program}\label{sec:first-program}
%%-----------------------
\cindex{first program}%

Let us consider the Goldstein-Price function
\begin{align*}
f(x,y)=&\bigl(1+(x+y+1)^2(19-14x+3x^2-14y+6xy+3y^2)\bigr)\\
&\bigl(30+(2x-3y)^2(18-32x+12x^2+48y-36xy+27y^2)\bigr)
\end{align*}
and let us enclose its range for $x\in\itv{-2}{2}$ and $y\in\itv{-2}{2}$.

\begin{Ex}\label{ex:goldstein-price}
The program below computes the formula of $f$ with intervals in place of
numbers. The line under it is what it prints.
 \begin{example}
 #include <iostream>
 #include <gaol/gaol.h>
 using namespace gaol;

 int main(void)
 {
   interval
     x(-2,2),
     y(-2,2), z;

   z=(1+sqr(x+y+1)*(19-14*x+3*sqr(x)-14*y+6*x*y+3*sqr(y)))*
     (30+sqr(2*x-3*y)*(18-32*x+12*sqr(x)+
        48*y-36*x*y+27*sqr(y)));

   std::cout << "z = " << z << std::endl;
   gaol::cleanup(); @rem^// Always, after the last use of GAOL~
   return 0;
 }
 @outputs^z = [-87881320, 147125080]~
 \end{example}
\end{Ex}

The program includes the header \file{gaol/gaol.h}, which gives access to
every facility of GAOL. The type \code{interval} and the functions of GAOL are
in the namespace \code{gaol\_core}, which a program names through the
namespace \code{gaol}: the \code{using} directive of Line~3 imports it. The
\file{gaol/gaol.h} header of GAOL~4 imported the namespace itself; that of
GAOL v5 no longer does, and a program adds the directive, or names
\code{gaol::interval} and \code{gaol::sqr}
(\pxref{sec:namespaces}).\newinvfive
\cindex{gaol!namespace}%
\cindex{namespace!gaol}%
\cindex{header!and namespace|see{namespace}}%

Nothing has to be called before using GAOL: \code{gaol::init()} runs by
itself when the program starts, before \code{main()}, with every compiler
(\pxref{chap:init}), and each operation of GAOL sets the rounding direction
it needs (\pxref{sec:rounding-direction}).\newinvfive

\begin{Rem}\label{rem:cleanup}
The program calls \code{gaol::cleanup()} right after its last use of GAOL, at
the latest at the end of \code{main()}, whatever the compiler. Unless GAOL is
built to preserve the rounding direction, its initialization sets the rounding
direction upward for the whole program, not only for the operations of GAOL,
and \code{gaol::cleanup()} sets back the one the program started with. GAOL
also calls it itself, but only when the program ends: without the call, the
code after the last use of GAOL would compute its doubles rounded
upward.\newinvfive
\end{Rem}

The function \code{sqr(\var{x})} is the square of \var{x}, and is equivalent
to \code{pow(\var{x},2)}: it is the function that \code{pow(\var{x},2)}
calls, with the same bounds, the sign of a zero bound
included.\newinvfive

\medskip
Let \file{f.cpp} be the name of the file containing the program above. It is
compiled with the flags that pkg-config gives (\pxref{sec:compiling}):
\cindex{compiling!a file using gaol}%

\begin{cmdshell}
% c++ -std=c++11 -O2 $(pkg-config --cflags gaol) -o f f.cpp \
      $(pkg-config --libs gaol)
\end{cmdshell}

\noindent
GAOL is the only library to link: CORE-MATH, which computes its elementary
functions, is compiled into it, and the C++ compiler links the mathematical
library of the system itself.

Executing \file{f} prints the line shown under the program. We then know that
$f(x,y)$ lies in \itv{-87881320}{147125080} for every $x$ and $y$ in
\itv{-2}{2}: \code{z} \emph{encloses} the range of $f$. It is not the range,
which is \itv{3}{1015690.27\ldots}: the minimum is $f(0,-1)=3$, and the
maximum is reached on the edge $y=2$. \code{z} is 231 times wider. Each
occurrence of $x$ and of $y$ in the formula ranges over \itv{-2}{2} on its
own, as if the others did not exist: this is the dependency problem
(\pxref{sec:dependency}), which makes an enclosure wider than the range.
\cindex{dependency problem}%
Cutting $x$ and $y$ into smaller intervals, and taking the hull of the values
of $f$ over the pieces, narrows the enclosure, at the price of more
evaluations; \file{examples/03\_dependency\_problem.cpp} of the sources does
it.

\section{The headers}\label{sec:headers}
%%-------------------
\cindex{headers}%

A program includes \file{gaol/gaol} (or \file{gaol/gaol.h}), and
\file{gaol/parser/gaol\_expression.h} for the expressions, with
\file{gaol/parser/gaol\_expr\_eval.h} to evaluate them. The three builds
install these headers and the headers they include. Each of the installed
headers compiles alone, but for the four that other headers include in the middle
of their code (\file{gaol/core/gaol\_interval\_fpu.h},
\file{gaol/core/gaol\_interval\_neon.h},
\file{gaol/core/gaol\_interval\_sse.h}) or that Visual C++ alone uses
(\file{gaol\_fpu\_msvc.h}). The headers of the sources of GAOL are not
installed:\newinvfive{} \file{gaol/parser/gaol\_interval\_parser.h}, which did
not compile when included, \file{gaol\_init\_cleanup.h},
\file{gaol/math/gaol\_core\_math.h}, \file{gaol/math/gaol\_u128.h},
\file{gaol/math/gaol\_compoundm1.h} and \file{gaol/math/gaol\_roundeven.h}.
Three headers of GAOL~4 are gone: \file{gaol\_parameters.h}, which declared
nothing, \file{gaol\_exact.h} with \file{sysdeps/}, which \file{dtoa.c}
needed, and \file{gaol\_assert.h}, whose \code{GAOL\_ASSERT} is
\code{assert()} now.

\medskip
The headers of GAOL compile without a warning under \option{-Wall -Wextra
-Wold-style-cast -Wundef}, and put nothing into the global namespace.
\file{gaol/gaol\_exceptions.h} declared \code{string} and \code{exception}
there; the header is gone, GAOL throws the exceptions of the standard library
(\pxref{sec:errors-overview}), and a program names \code{std::string} and
\code{std::exception}.\newinvfive{} The macros the headers define start with
\code{GAOL\_} (or \code{gaol\_}), so as not to meet those of the program or of
another library:\newinvfive{} \code{USING\_SSE2\_INSTRUCTIONS},
\code{HAVE\_FENV\_H}, \code{WORDS\_BIGENDIAN}, \code{INLINE},
\code{MEMALIGN}, \code{\_\_HI}\ldots{} are now
\code{GAOL\_USING\_SSE2\_INSTRUCTIONS}, \code{GAOL\_HAVE\_FENV\_H},
\code{GAOL\_WORDS\_BIGENDIAN}, \code{GAOL\_INLINE}, \code{GAOL\_MEMALIGN},
\code{GAOL\_HI}\ldots{} The header \file{<stdlib.h>} is no longer included
with \code{\_XOPEN\_SOURCE} defined again, and \code{nb\_fp\_numbers()}
returns an \code{unsigned long long}, which GAOL~4 named with the macro
\code{ULONGLONGINT}.

\medskip
The doubles \code{pi}, \code{half\_pi}, \code{two\_pi}, \code{pi\_dn},
\code{pi\_up}, \code{half\_pi\_dn}, \code{half\_pi\_up}, \code{ln2\_dn},
\code{ln2\_up}, \code{two\_power\_51} and \code{two\_power\_52} of GAOL~4 are
no longer declared for the program, which found them with \code{using
namespace gaol}, its own \code{pi} being then ambiguous.\newinvfive{} The
bounds of $\pi$ are those of \code{interval::pi()}, and the doubles around
$\pi$ and $\pi/2$ are \code{gaol\_pi\_dn}, \code{gaol\_pi\_up},
\code{gaol\_half\_pi\_dn} and \code{gaol\_half\_pi\_up}, in \code{gaol},
\code{gaol\_ieee1788} and \code{gaol\_core}.

\section{Compiling a program using GAOL}\label{sec:compiling}
%%--------------------------------------
\cindex{flags of interval arithmetic}%
\cindex{compiling!a file using gaol}%

The operations of GAOL on intervals are inline: the code that includes the
headers of GAOL has to be compiled with the \dfn{flags of interval
arithmetic}, not only GAOL itself.\newinvfive{} With GCC and Clang, each where
the compiler takes it, the flags are the following.

\begin{itemize}
\item \option{-frounding-math -fno-fast-math -ffp-contract=off}, so that the
  compiler rounds each operation in the direction in effect, as written, and
  neither evaluates it at compile time in the default rounding nor contracts a
  multiplication and an addition into a fused one.
\item \option{-msse2 -mfpmath=sse} on 32-bit x86, so that doubles are computed
  in double precision rather than on the x87 unit.
\item \option{-msse2 -msse3} on x86 processors with the SSE2 intervals
  (\option{GAOL\_SIMD}).
\item \option{-faligned-new=8} with GCC and Clang on 32-bit Windows
  (MinGW-w64, MSYS2) with the SSE2 intervals, so that the compiler asks for
  aligned memory for the intervals, which the C runtime would align on 8 bytes
  instead of 16.\newinvfive{} \code{gaol::gaol} and \file{gaol.pc} give it to
  the code using GAOL, which allocates the intervals, and
  \file{gaol/gaol\_config.h} checks it. It is an option of C++: GCC warns of a
  C file compiled with it (``valid for C++/ObjC++ but not for C''), and with
  \option{-Werror} the warning is an error. This happens where the flags of
  \file{gaol.pc}, or \code{gaol\_dep} of meson, are given to every file of a
  target, whereas \code{gaol::gaol} of CMake gives them to the C++ files only:
  a project that compiles C files gives the option to its C++ files only,
pkg-config having no flag for one
  language. Visual C++ and clang-cl have no such flag, and need none: their
  SSE2 intervals on 32-bit Windows keep their two bounds as two doubles, whose
  alignment is 8 (\code{GAOL\_SSE2\_UNALIGNED\_STORAGE} of
  \file{gaol/gaol\_config.h}).
\item \option{-mstackrealign} (or \option{-mincoming-stack-boundary=2}) with
  GCC on 32-bit x86, only for code that something calls with the stack aligned
  on 4 bytes, and not on 16 as GCC takes it (bug 40838 of GCC): a caller that
  does not follow the ABI of Linux, a hand-written trampoline or a foreign
  runtime for instance. A function of the program that keeps intervals in
  locals then stops on a \code{movaps} at an address that is not aligned
  (\code{test\_stack\_entered\_aligned\_on\_4\_bytes} of
  \file{tests/misc.cpp}, \pxref{sec:test-programs}).\newinvfive{} The
  programs for 32-bit Windows need neither: the system enters its threads and
  its callbacks with the stack aligned on 4 bytes, and that test passes
  without them with GCC for MinGW-w64, with Clang and with clang-cl, in the
  continuous integration.
\item Nothing more on 64-bit ARM processors with the NEON intervals
  (\option{GAOL\_SIMD}, \pxref{sec:neon}): the first item is all they ask
  for.\newinvfive
\item \option{-ffloat-store} where doubles are still computed on the x87 unit.
\item \option{-mfma} on x86, and \option{-mfpu=neon-vfpv4 -mfloat-abi=hard} on
  32-bit ARM, where GAOL is compiled with the fused multiply-add instructions
  of the processor (\option{GAOL\_FMA}, \pxref{sec:cmake-options}): the code
  using GAOL then runs only on a processor with them, as GAOL itself. With GCC
  for Windows on x86-64 and on 32-bit x86 the flags are \option{-mfma
  -mprefer-vector-width=128}, the second keeping the vectorizer of GCC to the
  128-bit registers, which it would otherwise spill to a stack that Windows
  aligns on 16 bytes, not 32 (GCC bug 54412): the objects of GAOL, CORE-MATH
  and the tests built so by GCC~16.2 have no instruction on a 256-bit
  register. A program that writes 256-bit intrinsics itself is exposed to that
  defect, as with \option{-mavx}, and so is one built with
  \option{-mavx512f} or a \option{-march} that turns on wider registers, which
  the option may not always undo. Cygwin and MSYS get these flags as the other
  GCC for Windows; Cygwin is built and tested by a job of the continuous
  integration, and refused without them.\newinvfive
\end{itemize}

\noindent
With Visual C++ and clang-cl, the flags are \option{/fp:strict},
\option{/arch:AVX2} for x64 with \option{GAOL\_FMA}, and \option{-msse2} with
clang-cl for 32-bit x86, which computes doubles on the x87 unit with
\option{/arch:SSE2}. Each build installs these flags with GAOL, in
\code{gaol::gaol} and \file{gaol.pc}. \file{gaol/gaol\_config.h} refuses code
compiled by Visual C++, or by clang-cl from Clang 16, without
\option{/fp:strict}. GCC and Clang do not tell the code whether
\option{-frounding-math} and \option{-ffp-contract=off} were given: there, it
only refuses what contradicts them, \option{-ffast-math},
\option{-ffinite-math-only} and doubles computed on the x87 unit
(\pxref{sec:refused}). The order of the options matters.
\option{-fno-fast-math} turns the first two off when it comes after them on
the command line, and does nothing when \option{-ffast-math} or
\option{-ffinite-math-only} comes after it, where the compilation stops.
\option{-Ofast}, which turns \option{-ffast-math} on, is refused with it,
except with GCC when \option{-fno-fast-math} is anywhere on the command line.
With Clang, \option{-ffast-math} and \option{-Ofast} also turn
\option{-frounding-math} off, which \option{-fno-fast-math} does not turn on
again: they have to come before the flags of interval arithmetic, not after
them.

\begin{Rem}\label{rem:gcc12-denorm}
With \option{-frounding-math} on x86, GCC 12.1 to 12.3 and 13.1 to 13.2
(Debian 12 has GCC 12.2) initialize wrongly a double member of a static array
of structures given \code{std::numeric\_limits<double>::denorm\_min()}: they
store the long double $2^{-1074}$ there, whose first eight bytes read $-0$ and
whose last two overwrite the next member (\code{\{ denorm\_min(), DBL\_MAX
\}} read $[-0, \mathrm{NaN}]$, the empty set, on 32-bit x86). A double or an
array of doubles initialized so is right, and so are GCC 12.4, 13.3 and
later. A program built with these releases computes such a value at run time,
as \file{tests/rounding\_direction.cpp} does.\newinvfive
\end{Rem}

The link of the program takes one option too, where the compiler accepts it
(GCC 13 and later on x86, and from 11.4 and 12.4 in the series 11 and 12):
\option{-mno-daz-ftz}, which keeps out of the link the file that sets the
flush-to-zero and denormals-are-zero modes when a program is linked with
\option{-Ofast} (\pxref{sec:ftz}). \code{gaol::gaol} and \file{gaol.pc} carry
it too, where the compiler that built GAOL accepts it: the three builds test
the compiler for it with a link.\newinvfive

\subsection{From CMake}\label{sec:compiling-cmake}
%%.....................
\cindex{CMake!using gaol from}%

A CMake project takes everything from the target \code{gaol::gaol} of the
package of GAOL:

\begin{example}
find_package(gaol REQUIRED)
target_link_libraries(my_target PRIVATE gaol::gaol)
\end{example}

\noindent
\code{gaol::gaol} carries the include directory, the flags above, and the
link option above for the code built by GCC~13 or a later one, or by the
compiler that built GAOL, of its major version and not older: another
compiler, as Clang~18, would stop on it, and so would GCC 12.0 to 12.3 with a
GAOL built by GCC 11.4, GCC having the option from 11.4 and 12.4 only in the
series 11 and 12. It no longer carries, for Visual C++, the definition
\code{\_\_GAOL\_PUBLIC\_\_=} of GAOL~4: GAOL being a static library there, its
declarations carry no \code{\_\_declspec(dllimport)}, and a definition left in
a project changes nothing.\newinvfive{} There is no other library to link:
CORE-MATH is compiled into \file{libgaol} itself. A library whose headers
include those of GAOL, as those of Codac do, links \code{gaol::gaol}
\code{PUBLIC}, so that its own users get the flags, and its CMake package
finds GAOL again (\code{find\_dependency(gaol)}). \file{tests/find\_package}
is a project using an installed GAOL this way.

A project can also build GAOL for itself, with the options it wants:

\begin{example}
include(FetchContent)
FetchContent_Declare(gaol GIT_REPOSITORY
    https://github.com/Jordan08/GAOL.git GIT_TAG master)
FetchContent_MakeAvailable(gaol)
target_link_libraries(my_target PUBLIC gaol::gaol)
\end{example}

\noindent
GAOL is then a target of the project, built with it, and \cmd{cmake --install}
of the project installs it with it, CMake package and \file{gaol.pc}
included; nothing is downloaded beyond the sources of GAOL.
\file{tests/fetch\_content} is a project building GAOL this way. GAOL is built
in Release where the project gives no build type, and
\option{-DGAOL\_DEBUG=ON} builds it for debugging whatever the build type or
the configuration of the project, whose own code keeps them
(\pxref{sec:cmake-options}).\newinvfive{} GAOL, a static library, is compiled
as position-independent code (\option{-fPIC}), so that \code{gaol::gaol} can
be linked into a shared library, such as Python bindings. A project that sets
\code{CMAKE\_POSITION\_INDEPENDENT\_CODE}, \code{ON} or \code{OFF}, before
\code{FetchContent\_MakeAvailable(gaol)}, or on the command line, is followed
instead.

\subsection{From pkg-config}\label{sec:compiling-pkgconfig}
%%..........................
\cindex{pkg-config}%

Each build installs \file{gaol.pc} in the \file{pkgconfig} directory of its
library directory (\file{\var{prefix}/lib/pkgconfig}, or \file{lib64} or
\file{lib/\var{multiarch}} rather than \file{lib} on the systems whose
libraries go there), except the CMake build with Visual C++. Its
\code{Cflags} carries the flags above with the include directory, and its
\code{Libs} GAOL itself and the link option above, where the compiler that
built GAOL accepts it. The variable \env{PKG\_CONFIG\_PATH} names the
directory to pkg-config:

\begin{cmdshell}
% export PKG_CONFIG_PATH=@var^prefix~/lib/pkgconfig
% c++ -std=c++11 -O2 $(pkg-config --cflags gaol) program.cpp \
      $(pkg-config --libs gaol)
\end{cmdshell}

\noindent
Built by CMake with its default options on x86-64 Linux, GAOL gives, on one
line:

\begin{cmdshell}
% pkg-config --cflags --libs gaol
-frounding-math -fno-fast-math -ffp-contract=off
-msse2 -msse3 -mfma -I@var^prefix~/include -L@var^prefix~/lib -lgaol
\end{cmdshell}

\noindent
The option \option{-mno-daz-ftz} comes after \option{-lgaol} where the
compiler that built GAOL takes it: \file{gaol.pc} is written for that
compiler, and a program linked by one that refuses the option, as Clang~18, a
GCC before 11.4 or GCC 12.0 to 12.3, stops on it, and is to be linked without
it. There is no \option{-lm}: GAOL still calls functions of the math library of
the system (those of \file{fenv.h}, \code{sqrt()}, \code{floor()},
\code{fma()} where the processor has no such instruction\ldots), which the
C++ compiler links itself.\newinvfive

\subsection{From meson}\label{sec:compiling-meson}
%%.....................
\cindex{meson!using gaol from}%

In a meson project, \code{dependency('gaol')} finds an installed GAOL. A
meson project can also build GAOL for itself, as a subproject, the sources of
GAOL being in its \file{subprojects/gaol}:\newinvfive

\begin{example}
gaol_dep = dependency('gaol', fallback: ['gaol', 'gaol_dep'])
executable('program', 'program.cpp', dependencies: gaol_dep)
\end{example}

\noindent
\code{dependency()} takes an installed GAOL where it finds one, and builds the
subproject otherwise; \code{subproject('gaol').get\_variable('gaol\_dep')}
always builds it. \code{gaol\_dep} carries the include directory, the flags
above and the link option above, as \file{gaol.pc} does: the flags that the
\file{meson.build} of GAOL gives its own targets do not reach those of the
project. \file{tests/meson\_subproject} is a project building GAOL this way.
GAOL could not be a subproject before GAOL v5: meson refuses the
\code{add\_global\_arguments()} its \file{meson.build} called in a subproject.

\medskip
GAOL is built in \code{release}, as CMake and \cmd{configure} build it brought
into another project, where the project leaves the build type to its default,
\code{debug}. meson 1.8.4 and later give a subproject its own default build
type, and a \loption{buildtype=debug} asked for builds GAOL for debugging.
Earlier ones give it the build type of the project, which would have built
GAOL without optimization, and before meson 1.8 with \code{GAOL\_DEBUGGING}:
there, where the options are those of \code{debug}, GAOL compiles its
libraries as in \code{release} all the same, as it cannot tell the default
from a \loption{buildtype=debug} asked for. \code{-Dgaol:enable-debug=true}
(or \code{default\_options: ['enable-debug=true']} in \code{subproject()} or
\code{dependency()}) builds GAOL for debugging, with any meson. Two cases keep
the build type of the project: Visual C++ with meson before 0.57, whose
\code{debug} gives \option{/RTC1}, which \cmd{cl} refuses with \option{/O2},
and the Visual Studio backend of meson (\loption{backend=vs}), which takes the
build type of the project for every target; the default backend, ninja, is not
one.\newinvfive

\section{The namespaces}\label{sec:namespaces}
%%----------------------
\cindex{namespace!gaol}%
\cindex{namespace!gaol\_core}%
\cindex{namespace!gaol\_ieee1788}%

The type \code{interval} of GAOL and its functions are in the namespace
\code{gaol\_core}. A program names them through one of two namespaces, which
\file{gaol/gaol.h} declares and does not open:\newinvfive

\begin{itemize}
\item[$\bullet$] \textbf{\code{gaol}}, GAOL under its own names, which this
  manual uses, to stay compatible with GAOL~4: a program written for GAOL~4
  keeps its names there and, but for the results that change
  (Chapter~\ref{chap:changes}), its results;
\item[$\bullet$] \textbf{\code{gaol\_ieee1788}}, GAOL under the names of IEEE
  1788-2015 (Chapter~\ref{chap:ieee1788}), to follow the standard: its
  operations have the names, the arguments and the results the standard gives
  them.
\end{itemize}

\noindent
The two namespaces hold the same type and compute with the same functions.
The main difference between them is the power, \code{pow} and \code{pown}
(\pxref{sec:pow-namespaces}), and the names that \code{textToInterval} reads.
A program that opens \code{gaol} writes the operations under their names:

\begin{example}
#include <gaol/gaol.h>
using namespace gaol;

interval x(1, 2);
interval y = pow(x, 3) + sin(x);
\end{example}

\noindent
\file{gaol/gaol.h} no longer opens \code{gaol} itself, as GAOL~4 did: a
program that relied on it adds \code{using namespace gaol;}, or names
\code{gaol::interval}.\newinvfive

\medskip
The two namespaces hold \code{interval}, and the functions that have the same
name and the same meaning in both (\code{sin}, \code{exp}, \code{sqrt},
\code{min}\ldots) are the same functions, those of \code{gaol\_core}.
\code{pow} and \code{textToInterval} are the ones whose meaning differs: each
namespace has its own, and none is in \code{gaol\_core}, the namespace where
argument-dependent lookup looks for a call \code{pow(x,~y)} on an interval, so
that each namespace finds its own only. A program opens one of the two
namespaces, not both: with both open, \code{pow(x,~y)} is ambiguous, and so is
\code{textToInterval(s)}, which reads the names of GAOL in \code{gaol} and
those of the standard in \code{gaol\_ieee1788} (\pxref{sec:constructors}).

\medskip
\code{gaol} names the type and the functions of \code{gaol\_core} one by
one.\newinvfive{} They are the names GAOL~4 had in \code{gaol} that GAOL v5
still names there (the class \code{interval}, the expressions with their nodes
and visitor, the functions and the operators of intervals, the functions of
doubles and the rounding direction, \code{round\_downward()},
\code{round\_upward()}, \code{round\_nearest()}, \code{previous\_float()},
\code{next\_float()}, \code{is\_finite()}, the version of the library,
\code{version\_major}, \code{version\_minor}, \code{version\_micro} and
\code{version}, and the doubles around $\pi$ and $\pi/2$ of
\exref{sec:headers}), the functions on intervals GAOL v5 adds (\code{exp2},
\code{log2}, \code{hypot}, \code{compoundm1}, \code{trunc}, \code{sinpi},
\code{erf}, \code{fma}, \code{cancel\_minus},
\code{round\_ties\_to\_even}\ldots{} and the nodes of their expressions),
\code{restore\_rounding()} and \code{exact\_string()}. \code{gaol\_ieee1788}
names the version as \code{gaol} does.

\noindent
Some names are in \code{gaol\_core} only, and a program names them there:
\code{gaol\_core::rnd\_keep}, for instance.\newinvfive{} They are
\code{rnd\_keep()} (GAOL v5 adds it, and had it in \code{gaol} in its 4.3
releases; GAOL~4 has none), \code{NaN\_val}, \code{eval\_stack}, and the
reader of strings, \code{parse\_interval()} and \code{parsing\_names}.

\noindent
What the code of GAOL uses for itself is in \code{gaol\_core} only too, out of
the way of the names of a program that opens \code{gaol}: the functions saving
and setting the rounding direction and the flush-to-zero modes
(\code{rounding\_state}, \code{get\_rounding()}, \code{set\_rounding()}, and
\code{round\_upward\_if\_needed()}, with \code{gaol\_set\_rounding\_x86()},
which GAOL v5 had in \code{gaol} in its 4.3 releases and GAOL~4 had not, and
\code{get\_flush\_modes()}\ldots), the functions computing the \code{pow} of
\code{gaol} (below), and the internals GAOL~4 had in \code{gaol} without its
manual documenting them: the functions \code{gaol\_signbit()},
\code{modulo\_k\_pi()} and the \code{round\_*\_sse()}, the types
\code{uintdouble}, \code{ullidouble}, \code{Interval\_struct} and
\code{prec\_t}, and the variable \code{the\_null\_expr}. A program calling one
of them names it in \code{gaol\_core}: \code{gaol\_core::get\_rounding()},
for instance.

\medskip
The \code{pow} of \code{gaol} takes the integer power for an integer exponent,
a negative base included, and the \code{pow} of IEEE 1788-2015 for any other
exponent (\pxref{sec:powers}). In \code{gaol\_core}, its versions are named
apart: \code{gaol\_pown()}, \code{gaol\_uipow()}, \code{gaol\_pown\_ll()},
\code{gaol\_pow\_real()}, \code{gaol\_pow\_hybrid()}, \code{gaol\_pown\_exp()}
and \code{gaol\_pow\_exp()}, which are not in \code{gaol}.

\section{The powers in the two namespaces}\label{sec:pow-namespaces}
%%----------------------------------------
\cindex{pow!in the two namespaces}%

The power is where \code{gaol} and \code{gaol\_ieee1788} differ.\newinvfive{}
IEEE 1788-2015 has two powers (Table 9.1): \ieee{pown}$(x,n)$, the integer
power, for an integer $n$, defined for every $x$, $x^0$ being $1$ and $0^n$
having no value for $n<0$; and \ieee{pow}$(x,y)$, the real power
$e^{y\log x}$, for an interval $y$, defined for $x>0$, and for $x=0$ when
$y>0$. GAOL~4 had a single \code{pow}, which took the power of a negative
base, and \code{gaol} keeps it: its \code{pow} is a hybrid of the two.

\begin{itemize}
\item In \code{gaol\_ieee1788}, \code{pown(x, n)} is the integer power, for an
  $n$ of any integer type, and \code{pow(x, y)} the real power: only the part
  of $x$ in \itv{0}{+\infty} counts, whatever $y$, and \code{pow(x, p)} with a
  number $p$ is \code{pow(x, [p])}, an integer $p$ included, as the standard
  gives no other meaning to it.
\item In \code{gaol}, \code{pow(x, y)} takes the integer power where the
  exponent is an integer: an integer type, a double that is an integer, or a
  degenerate interval \itv{n}{n} holding an integer, beyond the \code{int}s
  too; any other exponent takes the real power of the standard. There is no
  \code{pown} in \code{gaol}.
\end{itemize}

\begin{table}[htbp]\centering\renewcommand{\arraystretch}{1.1}
  \begin{tabular}{lll}
    \hline
    & \multicolumn{1}{c}{\code{gaol}} &
    \multicolumn{1}{c}{\code{gaol\_ieee1788}}\\
    \hline\hline
    \code{pow(\itv{-4}{-1}, 2)} &
    \itv{1}{16} &
    $\emptyset$ (\code{pown}: \itv{1}{16})\\
    \code{pow(\itv{-2}{3}, 3)} & \itv{-8}{27} & \itv{0}{27}\\
    \code{pow(\itv{0}{0}, 0)} &
    \itv{1}{1} &
    $\emptyset$ (\code{pown}: \itv{1}{1})\\
    \code{pow(\itv{-2}{2}, -1)} & \itv{-\infty}{+\infty} & \itv{0.5}{+\infty}\\
    \code{pow(\itv{-2}{-1}, \itv{10^{10}}{10^{10}})} &
    \itv{1}{+\infty} &
    $\emptyset$\\
    \code{pow(\itv{-1}{4}, 0.5)} & \itv{0}{2} & \itv{0}{2}\\
    \code{pow(\itv{0}{2}, -2)} & \itv{0.25}{+\infty} & \itv{0.25}{+\infty}\\
    \code{pow(\itv{-4}{-1}, \itv{1.5}{2.5})} & $\emptyset$ & $\emptyset$\\
    \hline
  \end{tabular}
  \caption{The power of \code{gaol} and of \code{gaol\_ieee1788} on some
    arguments}\label{tab:pow-namespaces}
\end{table}

Table~\ref{tab:pow-namespaces} shows the two on some arguments. They agree on
$x\subseteq\itv{0}{+\infty}$ and on every exponent that holds no integer
alone; they differ where $x$ has a negative part and the exponent is an
integer, and at $0^0$. The \code{pow} of \code{gaol} is thus not monotone for
the inclusion of its exponent: \code{pow(\itv{-4}{-1}, \itv{2}{2})} is
\itv{1}{16}, and \code{pow(\itv{-4}{-1}, \itv{1.5}{2.5})}, whose exponent
holds $2$, is empty; the \code{pow} of the standard is. Both give the tightest
bounds, but for the integer power: within the ints it comes from exact
products, and is one double beyond the tightest where it is within twice the
bound of the error of its products, never for the square; beyond the ints it
comes from the \code{pow} of CORE-MATH, and is the tightest for an exponent
that is a double and within three doubles of it otherwise, beyond $2^{53}$
(Chapter~\ref{chap:accuracy}). A degenerate integer exponent of
\code{gaol\_ieee1788::pow} takes the integer power too, with its bounds:
\code{pow([0x1.fffffffffffffp+1], [7])} has its lower bound one double below
the tightest.

\medskip
The exponent 2 is the square of the intervals: \code{pow(x, 2)},
\code{pow(x, [2])}, \code{gaol\_ieee1788::pown(x, 2)} and, on the part of $x$
in \itv{0}{+\infty}, \code{gaol\_ieee1788::pow(x, [2])} are \code{sqr(x)}
itself, with its bounds and the sign of its zero bounds, where the integer
power made its square from exact products, the same bounds in a longer time
and, for some intervals holding 0, a zero bound of the other
sign.\newinvfive{} The exception flags it raises are those of \code{sqr(x)}
too. For the empty set, \code{pow(x, 2)}, \code{pow(x, 2u)},
\code{pow(x, 2LL)}, \code{pow(x, 2.0)}, \code{pow(x, [2])},
\code{gaol\_ieee1788::pown(x, 2)} and \code{gaol\_ieee1788::pow(x, [2])}
return as \code{sqr} does, before any check, and leave the rounding direction
as the program had it, and the modes that flush the subnormal numbers to zero
as it had them too; so do the real and interval forms of \code{pow} for any
exponent, an empty exponent included, and the \code{long long} form beyond
the unsigned ints.\newinvfive{} The integer powers of an \code{int} or of an
\code{unsigned int} other than 2 check the direction first and leave it
upward: a program does not count on the direction after an operation that
computes no bound (\pxref{sec:rounding-direction}).

\medskip
The rest of GAOL follows the same split. The expressions \code{pow(e, n)}
and \code{pow(e1, e2)} of \code{gaol} are the hybrid power, \code{pown(e, n)}
and \code{pow(e1, e2)} of \code{gaol\_ieee1788} the two powers of the
standard. The text \code{"pow([-4,-1],2)"} is read as \itv{1}{16} by the
\code{textToInterval} of \code{gaol}, and as the empty set by that of
\code{gaol\_ieee1788}, which reads \code{pow} as the real power and
\code{"pown([-4,-1],2)"} as the integer power, a name \code{gaol} does not
read. The reverse functions are those of the standard in both namespaces:
\code{pow\_rel} and \code{pow\_exponent\_rel} (\ieee{powRev1} and
\ieee{powRev2}) reverse the real power, and \code{nth\_root\_rel}
(\ieee{pownRev}) the integer power, for every integer exponent, where it took
$n\geq 1$ only.\newinvfive

A program written for GAOL~4, or for IBEX and Codac, whose intervals take the
\code{pow} of GAOL, opens \code{gaol}; a program that wants the results of the
standard, or compares GAOL with other libraries of IEEE 1788-2015, as the
comparison of \exref{sec:comparison} does, opens \code{gaol\_ieee1788}.

\begin{Ex}\label{ex:pow-namespaces}
The three powers of \itv{-4}{-1} with the exponent 2, named in their
namespaces.
\begin{example}
#include <iostream>
#include <gaol/gaol.h>

int main(void)
{
  gaol::interval x(-4, -1);
  std::cout << gaol::pow(x, 2) << ' '
            << gaol_ieee1788::pow(x, 2) << ' '
            << gaol_ieee1788::pown(x, 2) << std::endl;
  gaol::cleanup();
  return 0;
}
@outputs^[1, 16] [empty] [1, 16]~
\end{example}
\end{Ex}

\section{Integers}\label{sec:integers}
%%---------------
\cindex{integers!given to GAOL}%

An integer given to GAOL is taken as the integer it is.\newinvfive{}
\code{interval(n)} is the tightest interval containing $n$: \itv{n}{n} where
$n$ is a double, as every \code{int} is, and the two doubles around it
otherwise. \code{interval(9007199254740993LL)}, $2^{53}+1$, is
\itv{2^{53}}{2^{53}+2}, where its conversion to a double gave
\itv{2^{53}}{2^{53}} or \itv{2^{53}+2}{2^{53}+2}, according to the rounding
direction, neither of which contains it, and \code{interval(ULLONG\_MAX)}
contains $2^{64}-1$. The same holds for \code{interval(a, b)} with an integer
bound, which is empty for $a>b$ compared as numbers
(\code{interval(9007199254740993LL, 9007199254740992LL)} is empty, though both
lie between the same two doubles), for the assignment \code{x = n} and
\code{x |= n}, the arithmetic operators (\code{x + n}, \code{n * x},
\code{x += n}\ldots), \code{set\_contains(n)}, \code{expression(n)}, and
\code{numsToInterval()} and \code{isMember()} in \code{gaol\_ieee1788};
\code{x \&= n} is \code{interval(n)} where \code{x} contains $n$, and the empty
set otherwise. The relations and \code{set\_strictly\_contains(n)} compare $n$
as the integer it is: \code{interval(9007199254740992.0) < 9007199254740993LL},
$\itv{2^{53}}{2^{53}}<2^{53}+1$, is \code{true}, though no double lies between
them. The constructors are detailed in \exref{sec:constructors}, and the
relations in \exref{sec:relational-symbols}.

\medskip
Every integer type is taken so, \code{bool} aside, whatever the rounding
direction; an unscoped enumeration, and \code{\_\_int128} outside the GNU
modes of GCC and Clang, still convert to a double, and a \code{long double}
bound next to an integer takes \code{interval(double, double)}. An integer
that is a double gives the result the double gives, and for the types of 53
bits or less (\code{int}, \code{short}, \code{char}\ldots) the same code.
\code{interval(n)} and \code{expression(n)} are explicit, as their
constructors from a double; \code{expression(n)} is \code{expression(double)}
where $n$ is a double.

\medskip
\code{pow(x, n)} of \code{gaol} and \code{pown(x, n)} of \code{gaol\_ieee1788}
take an exponent of any integer type: \code{pow(x, 5L)} did not compile, a
\code{long} being as far from an \code{int} as from an \code{unsigned} and a
\code{double}, and \code{pown(x, 3000000000u)} converted the exponent to a
negative \code{int}. An exponent within the ints is taken as an \code{int},
one within the unsigned ints as an \code{unsigned}, and one beyond them as a
\code{long long} (\code{gaol\_pown\_ll()}), whose power comes from the
\code{pow} of CORE-MATH at the bounds of $x$, as for \code{gaol::pow(x, [n])}
(\itv{-\infty}{+\infty} for a nonempty $x$ before GAOL v5):
\code{pow(interval(1 + 0x1p-30), 1LL << 33)}, about $e^8$, is
\itv{2980.9579759367944}{2980.9579759367948}. Beyond 64 bits (an
\code{\_\_int128}), the exponent is taken as $\pm(2^{63}-1)$ or
$\pm(2^{63}-2)$, of its sign and parity, the power of every double but $0$,
$\pm 1$ and $\pm\infty$ being beyond the doubles there. The same holds for
\code{nth\_root(x, q)} and \code{rootn(x, q)}, whose order beyond the ints and
the unsigned ints gives an enclosure of the roots, computed by \code{pow} with
the exponent $1/q$ ($x^{1/q}$ is within $2^{-21}$ of $1$ for every double
$x>0$ there), for \code{pownRev(c, x, p)}, which gives the hull of the parts
of $x$ in these roots and, for an even $p$, in their opposites (for a negative
$p$, $c$ is cut at 0 first), and for \code{pow(e, n)}, \code{pown(e, n)} and
\code{nth\_root(e, n)} of the expressions. \code{nth\_root(x, 3L)} was
ambiguous; \code{rootn(x, 3000000000u)} and \code{pownRev(c, 3000000000u)}
converted the order to a negative \code{int}, the first giving a wrong
interval (\itv{0.99999999839}{0.99999999946} for \itv{2}{8}, rather than
\itv{1.00000000023}{1.00000000069}) and the second the reverse of a negative
power; \code{rootn(x, 5000000001LL)} and the powers of expressions cut a
\code{long} to an \code{int}; and \code{nth\_root(e, -2)} took the order
4294967294. Beyond the ints, \code{pow(e, n)} and \code{pown(e, n)} of the
expressions are the \code{pow} of GAOL with the exponent \itv{n}{n}, the
powers of the intervals where $n$ is a double, and beyond $2^{53}$ the
product of $x^h$ and $x^l$, $n=h+l$, $h$ a double and $|l|<2^{11}$, no node
holding a wider integer: wider than the power of the interval where $e$
contains 0; beyond the unsigned ints, \code{nth\_root(e, n)} is the root of
the interval but where $e$ contains 0, where it is wider.

\section{A result thrown away}\label{sec:nodiscard}
%%---------------------------
\cindex{result thrown away}\findex{GAOL\_NODISCARD}%

The functions of GAOL do not change the interval they are given: they return
another one. \code{sqrt(x);} alone leaves \code{x} as it was, and so do
\code{x.mid();} and \code{x.emptyset();}, a static function that returns the
empty set. The functions whose result is all they do carry an attribute that
makes the compiler warn about such a call:\newinvfive{} the functions and
operators on intervals, their predicates, the names of \code{gaol\_ieee1788}
and the functions building expressions.

\begin{example}
sqrt(x);        @rem^// warning: the result is thrown away, x is unchanged~
x = sqrt(x);    @rem^// what was meant~
(void)sqrt(x);  @rem^// thrown away on purpose: no warning, but with GCC before 7~
\end{example}

\noindent
The attribute is \code{GAOL\_NODISCARD}, of \file{gaol/gaol\_config.h}. It is
\code{[[nodiscard]]} in C++17 and later. A program is compiled in the standard
its project sets, or else in the one of its compiler, which is C++14 for GCC 6
to 10, Clang 6 to 15 and Visual C++. Before C++17, \code{GAOL\_NODISCARD} is
\code{[[nodiscard]]} too where the compiler takes it there: GCC 7 and later,
and Visual C++ 2019 16.4 and later, whose warning is C4834. Clang, which warns
about \code{[[nodiscard]]} before C++17 with \option{-pedantic}, and GCC
before 7 take \code{\_\_attribute\_\_((warn\_unused\_result))}, whose warning
is \option{-Wunused-result} (on by default); older Visual C++ takes
\code{\_Check\_return\_} of \file{<sal.h>}, which only the code analysis
reports (\option{/analyze}, warning C6031), not the compiler. A cast to void
(\code{(void)sqrt(x);}) silences each of them, but the attribute of GCC
before 7: the result is kept in a variable there. A program that does not
want the warnings defines \code{GAOL\_NODISCARD} empty before including GAOL
(\option{-DGAOL\_NODISCARD=}). The compound assignments (\code{x += y}), which
do change \code{x}, do not carry it.

\section{Tools of interval algorithms}\label{sec:tools}
%%-----------------------------------
\cindex{tools of interval algorithms}%
\findex{width\_enclosure}\findex{inflate}\findex{midrad}\findex{bisect}%
\findex{is\_bisectable}%

The algorithms of interval computing, the enclosure of a range by
subdivision, the interval Newton method, the contractors, the branch and
bound, are written with a few operations beyond the arithmetic, which they
otherwise write for themselves. GAOL v5 provides them.\newinvfive

\begin{itemize}
\item[$\bullet$] \textbf{The width.} \code{x.width\_enclosure()} is the
  interval \code{interval(r) - interval(l)} enclosing the exact width $r-l$ of
  $x=\itv{l}{r}$, whose upper bound is \code{width()}, the width rounded
  upward; it is \itv{\mathrm{MAX}}{+\infty} for an unbounded $x$, as
  \code{textToInterval()} reads \code{inf}, and the empty set for the empty
  set.
\item[$\bullet$] \textbf{The inflation.} \code{x.inflate(r)} is $x$ widened by
  the absolute radius $r$, \code{x + interval(-r, r)}, each bound rounded
  outward; $r=+\infty$ gives the whole line, and $r<0$ or a NaN $r$ the empty
  set.
\item[$\bullet$] \textbf{A midpoint and a radius.} \code{interval::midrad(m, r)}
  is the interval of midpoint $m$ and radius $r$,
  \code{interval(m) + interval(-r, r)}, which encloses $[m-r,\,m+r]$, where
  \code{interval(m - r, m + r)} does not: both bounds are then computed in the
  upward rounding of GAOL, the lower one above $m-r$ where it is inexact. An
  infinite $m$ gives the empty set, as \code{interval(m)} does.
\item[$\bullet$] \textbf{The bisection.} \code{x.bisect(ratio)} is the pair of
  the parts of $x$ on either side of a cut point at the given ratio of its
  width, $0<\mathit{ratio}<1$, which share the cut point and cover $x$;
  \code{bisect(0.5)} cuts where \code{split()} does, at \code{midpoint()}, and
  so does every ratio for a half-line or the whole line. A ratio outside
  $(0,1)$, or NaN, throws \code{std::domain\_error}. \code{x.is\_bisectable()}
  tells whether a double lies strictly between the bounds, both parts being
  then smaller than $x$: not for two consecutive doubles, nor for
  \itv{\mathrm{MAX}}{+\infty} (\pxref{sec:splitting}).
\item[$\bullet$] \textbf{The hull and the intersection.} \code{hull(x, y)} and
  \code{intersect(x, y)} are \code{x | y} and \code{x \& y} as functions, in
  \code{gaol} as in \code{gaol\_core}.
\end{itemize}

\begin{Ex}\label{ex:tools}
The program below uses these tools on an enclosure of $0.1$, on the interval
of midpoint $1$ and radius $10^{-20}$, and on \itv{0}{1} cut at a quarter of
its width.
\begin{example}
#include <iostream>
#include <gaol/gaol.h>
using namespace gaol;

int main()
{
  const interval x = textToInterval("0.1");        @rem^// encloses 0.1~
  const interval y = interval::midrad(1.0, 1e-20); @rem^// [1 - 1e-20, 1 + 1e-20]~
  const std::pair<interval, interval> parts = interval(0.0, 1.0).bisect(0.25);
  std::cout << x.width_enclosure() << ' ' << y << ' ' << parts.first << ' '
            << parts.second << ' ' << hull(x, y) << std::endl;
  gaol::cleanup();
  return 0;
}
\end{example}
\end{Ex}

\paragraph{The reverse functions.}
A contractor shrinks a box without losing a solution of a constraint: it
evaluates the expression forward, from the variables to the result, and goes
backward, each node being intersected with what the others allow. The
backward operations are the relational functions of GAOL, the \code{*\_rel}
functions, which reverse each operation (\pxref{sec:relational-arithmetic}).
Their arguments are in one order, the result first and the operand contracted
last: \code{div\_rel(c, b, x)} contracts a factor of $z=x\cdot y$, and
\code{sqrt\_rel(c, x)} the argument of a square root. The reverse functions of
the power follow it: \code{pow\_rel(K, J, I)} is the hull of the $x$ of $I$
with an $x^y$ in $K$ for a $y$ of $J$, and \code{pow\_exponent\_rel(K, H, I)}
the hull of the $y$ of $I$ with an $x^y$ in $K$ for an $x$ of $H$; they are
\ieee{powRev1}$(b,c,x)$ and \ieee{powRev2}$(a,c,x)$ of IEEE 1788-2015 (Table
10.1), as \code{pow\_rel(c, b, x)} and \code{pow\_exponent\_rel(c, a, x)}, $x$
being \itv{-\infty}{+\infty} when left out.\newinvfive{} Both reverse the
\code{pow} of the standard, that of \code{gaol\_ieee1788::pow}, in either
namespace: $x^y$ is $e^{y\log x}$ for $x>0$, $0$ for $x=0$ and $y>0$, and has
no value elsewhere, so that no $x<0$ is a solution, and
\code{pow\_rel([1, 4], [2], [-3, 3])} is \itv{1}{2}. The \code{pow} of
\code{gaol} takes the integer power for a degenerate integer exponent
\itv{n}{n}, which has negative bases too: it is not one function of $x$ and
$y$, and the reverse of $x^n$ is \code{nth\_root\_rel(K, n, I)}, which takes
every integer $n$ (\pxref{sec:relational-arithmetic}).

\medskip
\code{max\_rel(z, y, x)}, \code{min\_rel(z, y, x)}, \code{sign\_rel(z, x)} and
\code{floor\_rel(z, x)} contract $x$ under the constraints $z=\max(x,y)$,
$z=\min(x,y)$, $z=\mathrm{sign}(x)$ and $z=\lfloor x\rfloor$.\newinvfive{}
Each returns the hull of the $x$ of \code{x} that the function sends into
\code{z}, with some $y$ of \code{y} for \code{max} and \code{min}: the reverse
functions IEEE 1788-2015 defines for every function (10.5.4), with the
arguments in the order of the relational functions of GAOL. \code{max} and
\code{min} being symmetric, \code{max\_rel(z, x, y)} contracts $y$, and
\code{max(x, y) \& z} contracts $z$. IEEE 1788-2015 names and requires only
the reverse functions of its Table 10.1, which these four are not in:
\code{gaol\_ieee1788} has no name for them, and a program that opens it calls
them as \code{max\_rel(z, y, x)} all the same, argument-dependent lookup
finding them in \code{gaol\_core} (\pxref{sec:namespaces}).

\begin{example}
interval x(-5.0, 5.0), z(1.0, 2.0);
interval a = max_rel(z, interval(0.0, 1.0), x);   @rem^// [-5, 2]: y meets z~
interval b = max_rel(z, interval(-1.0, 0.0), x);  @rem^// [1, 2]: y below z~
interval c = max_rel(z, interval(3.0, 4.0), x);   @rem^// empty: y above z~
interval s = sign_rel(interval(-1.0), x);         @rem^// [-5, 0]~
interval f = floor_rel(interval(0.0, 2.0), x);    @rem^// [0, 3]~
\end{example}

\medskip
The interval Newton method encloses the roots of $f$ in $X$ with the
relational division, which solves $(x-m)\,d=-f(m)$ for $d$ in the enclosure of
$f'(X)$ where the division \code{/} does not (\exref{sec:common-errors}, and
\file{examples/05\_interval\_newton.cpp}); dividing by the negative and by the
positive part of $f'(X)$ apart, two \code{div\_rel}, or
\code{mulRevToPair(b, c)}, keeps the gap around $m$. The examples 06 and 07
of \file{examples/} write a branch and bound and a contractor this way, with
\code{bisect()}, \code{is\_bisectable()} and the hull.

\section{Errors, in short}\label{sec:errors-overview}
%%-----------------------
\cindex{errors!overview}%

GAOL reports an error by throwing an exception of the standard library
(\file{<stdexcept>}), where GAOL~4 had classes of its own:\newinvfive{}
\code{std::invalid\_argument} for a string that is no interval, given to
\code{gaol::textToInterval()} or read by \code{operator>{}>}, and
\code{std::domain\_error} for a function called with an argument it does not
take, \code{nb\_fp\_numbers()} with a NaN, \code{bisect()} with a ratio
outside $(0,1)$, or \code{nth\_root(8, 1.5)} in a string. Both derive from
\code{std::logic\_error}, which derives from \code{std::exception}, so that a
handler that knows nothing of GAOL says what went wrong: its \code{what()} is
\code{\var{file}:\var{line}: \var{explanation}}, the name of the source file
of GAOL without its directories and the line that threw, then what went
wrong. Where memory runs out, GAOL throws the \code{std::bad\_alloc} of C++,
the reader of strings included. A build without
\code{GAOL\_EXCEPTIONS\_ENABLED} prints a message and aborts instead of
throwing; the three builds can be asked to do so (a setting of the developers). The classes of GAOL~4, what
became of each, the messages, the warnings and the builds that abort are
detailed in \exref{sec:exceptions}.

\section{Common errors}\label{sec:common-errors}
%%---------------------
\cindex{common errors}%

This section reviews the errors most often made with GAOL, so that the user
recognizes them in a program. Most of them come from the two facts that a
double is not an interval, and that an interval is not a number.

\subsection{Floating-point arithmetic and rounding}
%%.................................................
\cindex{Rounding!to nearest}%

Programming with floating-point numbers is one of the few activities where
the compiler is to be watched. Let us consider the following piece of code:

\begin{example}
#include <gaol/gaol>  @rem^// We do not import the gaol namespace~

using gaol::interval;

int main(void) {
  interval one_tenth(0.1); @rem^BEWARE: wrong !~

  @emph^[some code using one_tenth]~
  gaol::cleanup();
}
\end{example}

\noindent
Though a rational perfectly representable in decimal, \code{0.1} is not
representable in binary with a finite number of bits, and thus requires
rounding. The purpose of the user was to define an interval containing this
value. However, \code{0.1} is rounded \emph{at compile time}, most certainly
to the nearest representable floating-point number \roundNearest{0.1}
(\pxref{sec:rounding-definitions}). As a consequence, \code{one\_tenth} is a
degenerate interval containing only \roundNearest{0.1}, and the containment
property (\pxref{sec:containment}) is violated.

The right way to deal with rational constants that might not be perfectly
representable as floating-point numbers is to write them as text, and to read
them with \code{gaol::textToInterval()} (\pxref{sec:constructors}), so that
they are correctly rounded downward and upward at \emph{run time}:

\begin{example}
#include <gaol/gaol>

using gaol::interval;

int main(void)
{
  @rem^// OK: this is the right way~
  interval one_tenth = gaol::textToInterval("0.1");
  @rem^// Another possible way~
  interval one_tenth2 = gaol::textToInterval("1/10");

  @emph^[some code using one_tenth]~
  gaol::cleanup();
  return 0;
}
\end{example}

\noindent
Now, \code{one\_tenth} is the smallest floating-point interval enclosing
$0.1$. An interval like this one, containing at most two consecutive
floating-point numbers, is a canonical interval
(\pxref{sec:intervals}).
\cindex{interval!canonical|see{canonical interval}}%
\cindex{canonical interval}%
Written in hexadecimal (\pxref{sec:output-format}), which shows the bounds
exactly, the two intervals are:

\begin{onscreen}
interval(0.1) = [0x1.999999999999ap-4]
one_tenth     = [0x1.9999999999999p-4, 0x1.999999999999ap-4]
\end{onscreen}

\subsection{Functions of numbers}
%%...............................
\cindex{sqrt@\code{sqrt} of a number}%

A function called with a number is the function of the C library, even where
the namespace \code{gaol} is open: \code{sqrt(2.0)} is a \code{double},
rounded in the direction in effect, and no enclosure of $\sqrt{2}$. A
\code{double} no longer becomes an interval by itself, the constructor being
explicit, so that \code{interval r = sqrt(2.0);} does not compile; but
\code{interval r(sqrt(2.0));} does, and gives a point that does not contain
$\sqrt{2}$. The enclosure is the function of GAOL on an interval, or the text
read by \code{textToInterval()}:\newinvfive

\begin{example}
interval root2 = textToInterval("sqrt(2)");
cout << boolalpha << interval(sqrt(2.0)).set_contains(root2) << ' '
     << sqrt(interval(2.0)).set_contains(root2);
@outputs^false true~
\end{example}

\noindent
For the same reason, \code{gaol::sin(0.5)} does not compile: it is written
\code{gaol::sin(interval(0.5))}, and \code{sin(0.5)} alone is the \code{sin}
of C. Writing code inside \code{namespace gaol \{\}} has the same effect: the
using-declarations of \code{gaol} are visible there, a call \code{sqrt(2.0)}
finds \code{gaol::sqrt}, which hides \code{::sqrt}, and does not compile. This
was already the case with GAOL~4.

\subsection{Midpoint and radius}
%%..............................
\cindex{interval!midpoint and radius}%

\code{interval(m - r, m + r)}, with doubles, is not an interval enclosing
$[m-r,\,m+r]$: the program computes \code{m - r} and \code{m + r} itself, in
the direction in effect, which is upward until \code{cleanup()}, so that
\code{m - r} may be above the real $m-r$, and to nearest each bound may be on
either side. With \code{m = 1.0} and \code{r = 1e-20}, the interval is
\itv{1}{1+2^{-52}} while the direction is upward, and does not contain
$1-10^{-20}$. \code{interval::midrad(m, r)} (\pxref{sec:constructors}),
or \code{m + interval(-r, r)}, computes both bounds in intervals, and gives
\itv{1-2^{-53}}{1+2^{-52}}, which does.\newinvfive

\begin{example}
double m = 1.0, r = 1e-20;
interval wrong(m - r, m + r);  @rem^// Does not enclose [m - r, m + r]~
interval right = interval::midrad(m, r);  @rem^// Does~
\end{example}

\subsection{The doubles of the program}
%%.....................................

The rounding direction that GAOL leaves applies to the rest of the program,
which computes its own doubles in it: after an operation of GAOL,
\code{1.0 / 3.0} is the double above $1/3$, and \code{printf()} and the
streams of the GNU C library write the decimal digits of a double rounded
upward too, until \code{cleanup()}. \exref{sec:upward-effects} gives what
changes, and what to do.\newinvfive

\subsection{Zero}
%%...............
\cindex{interval!zero}%

\code{interval~x(0,~0)}, \code{interval~x(0)} and \code{interval::zero()}
are \itv{0}{0}. The literal \code{0} is also a null pointer, and the
constructors from one and from two strings of GAOL~4 made these calls
ambiguous; GAOL v5 has no constructor from strings.\newinvfive

\subsection{The sign of a zero bound}\label{sec:zero-sign}
%%...................................
\cindex{zero!sign of a bound}%

A bound that is zero has a sign, which differs between the builds:
\code{sqr(interval(-1, 2))} is \itv{-0}{4} with the SSE2 intervals, which the
default build uses on x86 processors, and with the NEON intervals, which it
uses on 64-bit ARM processors, and \itv{0}{4} with the FPU intervals, those of
the other processors and of the builds with \option{GAOL\_SIMD} off (Visual
C++ and clang-cl have the SSE2 intervals on x86 and on x64);\newinvfive{} \code{interval::zero()} is \itv{-0}{0} with the first
two and \itv{0}{0} with the third. \code{pow(x, 2)} is \code{sqr(x)}, and has
the same sign.\newinvfive{} The NEON and the SSE2 intervals can differ too:
the maximum of 64-bit ARM, which the product takes of its four products, ranks
$+0$ above $-0$, where the one of x86 returns its second operand, so that the
lower bound of \code{interval(-1, 0) * interval(-1, 0)} is $-0$ with the NEON
intervals and $+0$ with the others (\pxref{sec:neon}); the sign of a zero
bound of a product of the FPU intervals is the one the compiler gives
(\pxref{sec:fpu-intervals}). IEEE 1788-2015 and GAOL take the two zeros for
the same number: \code{x.left() == 0.0} is \code{true} whichever it is, and a
point interval at zero is written \code{[0]}. A program that reads the sign,
with \code{std::signbit(x.left())} or \code{1.0 / x.left()}, which is
$-\infty$ or $+\infty$, gets one answer with a build and the other with the
other: it compares the bounds with \code{==}, \code{<} and
\code{<=}.\newinvfive

\subsection{The empty set and the domain of a function}
%%.....................................................
\cindex{empty set!in tests}%
\cindex{domain of a function}%

A function is computed on the part of its argument that lies in its domain
(IEEE 1788-2015), and GAOL, which has no decorations, says nothing of the
rest: \code{sqrt(interval(-1, 4))} is \itv{0}{2}, \code{log(interval(-1, 0))}
is empty, and so are \code{interval(1, 0)} and the interval of one infinite
or NaN double. The empty set then passes the tests as if it were a good
answer: it has no point that contradicts a certainly relation, and it is
included in every interval, so that a box outside the domain of a function is
both inside and outside the set a paving looks for:

\begin{example}
interval e = sqrt(interval(-2.0, -1.0)), z(0.0, 1.0);
cout << boolalpha << e.is_empty() << ' ' << (e < z) << ' '
     << (e > z) << ' ' << z.set_contains(e) << ' '
     << e.set_disjoint(z);
@outputs^true true true true true~
\end{example}

\noindent
A program tests \code{is\_empty()} first, and decides itself what a box
outside the domain of its function means, as \file{examples/08\_sivia.cpp}
does.\newinvfive

\subsection{Relations, sorting and containers}
%%............................................
\cindex{relation!symbols!pitfalls}%

The symbols \code{<}, \code{<=}, \code{>} and \code{>=} are the certainly
relations of IEEE 1788-2015, which do not tell which of two intervals is the
greater, and there is no \code{==}: \code{interval(1,3) < interval(2,4)} and
\code{interval(1,3) >= interval(2,4)} are both \code{false}
(\pxref{sec:relational-symbols}). In C-XSC and in PROFIL/BIAS, \code{x <= y}
means that \code{x} is a subset of \code{y}, and \code{r <= x} that the number
\code{r} is in \code{x}: code ported from them compiles, and changes its
meaning (\pxref{sec:other-libraries}). None of them is an ordering for the
standard library: \code{std::sort}, \code{std::set} and \code{std::max} take
the comparator \code{gaol::lexicographic\_less}
(\pxref{sec:lexicographic-less}).\newinvfive

\subsection{Both namespaces}
%%..........................

A program opens the namespace \code{gaol} or the namespace
\code{gaol\_ieee1788}, not both (\pxref{sec:namespaces}):

\begin{example}
#include <gaol/gaol.h>
using namespace gaol;
using namespace gaol_ieee1788;

gaol::interval x(1,2), y = pow(x,x); @rem^// Does not compile~
\end{example}

\noindent
The compiler reports that the call of \code{pow} is ambiguous. Not every call
is: with an \code{int} or an \code{unsigned int} for exponent, \code{pow(x, 2)}
compiles, and takes the \code{pow} of \code{gaol}, silently, where that of
\code{gaol\_ieee1788} gives another result for a negative base
(\pxref{sec:pow-namespaces}). Only the calls with another type of exponent,
\code{pow(x, 2.0)}, \code{pow(x, 2L)} or \code{pow(x, y)}, are
ambiguous:\newinvfive

\begin{example}
#include <iostream>
#include <gaol/gaol.h>
using namespace gaol;
using namespace gaol_ieee1788;

int main(void)
{
  gaol::interval x(-4,-1);
  std::cout << pow(x, 2) << std::endl; @rem^// The pow of gaol~
  gaol::cleanup();
  return 0;
}
@outputs^[1, 16]~
\end{example}

\subsection{Bisection}
%%....................
\cindex{bisection!of an unbounded interval}%

A loop that cuts an interval until it is narrow,
\code{while (x.width() > eps)}, never ends on an unbounded interval:
\code{bisect(0.5)}, like \code{split()}, cuts \itv{1}{+\infty} at
\code{DBL\_MAX}, and the half \itv{\mathrm{DBL\_MAX}}{+\infty} is cut at
\code{DBL\_MAX} again, its width staying infinite. What a loop stops on is
\code{is\_bisectable()}, which is \code{false} for this half, and for an
interval of two consecutive doubles, the cut of which gives the interval
itself as one half (\pxref{sec:splitting}):\newinvfive

\begin{example}
interval x(1.0, GAOL_INFINITY);
x = x.bisect(0.5).second;
cout << boolalpha << x << ' ' << x.width() << ' '
     << x.is_bisectable();
@outputs^[1.797693134862315e+308, inf] inf false~
\end{example}

\subsection{Newton's method: the division and the relational division}
%%....................................................................
\cindex{Newton's method!division}%

The step $N(X)=m-f(m)/F'(X)$ of the interval Newton method divides with
\code{/}, the set of the quotients $a/b$ of a point of the numerator by a point
other than 0 of the denominator. Where $0\in F'(X)$ and $f(m)=0$, the quotient
$0/F'(X)$ is $\itv{0}{0}$, so that $N(X)=m$ lies inside $X$ and ``proves''
that $X$ holds exactly one root, where it holds seven for $\sin x-x/10$ on
\itv{-10}{10}. A root satisfies a relation, $(x-m)\,d=-f(m)$ for some $d$ in
$F'(X)$, which \code{\%}, or \code{div\_rel()}, the reverse of the product
(\pxref{sec:relational-arithmetic}), solves: it gives the whole line, rightly,
when $d$ and $f(m)$ can both be 0.

\begin{example}
interval fm(0.0), D(-1.1, 0.9);
cout << fm / D << ' ' << fm % D << ' '
     << div_rel(fm, D, interval(-5.0, 5.0));
@outputs^[0] [-inf, inf] [-5, 5]~
\end{example}

\noindent
The functional division is the right one where $F'(X)$ holds no 0.
\file{examples/05\_interval\_newton.cpp} shows the trap and the extended step.

\subsection{Reading text}
%%.......................
\cindex{textToInterval@\code{textToInterval}!errors}%

\code{gaol::textToInterval()} throws \code{std::invalid\_argument} for a text
that is no interval, and \code{std::domain\_error} for a text that is one but
calls a function with an argument it does not take, such as
\code{nth\_root(8, 1.5)}, whose order is not an integer: a loop that catches
the first lets the second escape. Both derive from \code{std::logic\_error}
(Chapter~\ref{chap:errors}).\newinvfive

\begin{example}
try {
  cout << textToInterval("nth_root(8, 1.5)");
} catch (const std::invalid_argument&) {
  cout << "std::invalid_argument";
} catch (const std::domain_error&) {
  cout << "std::domain_error";
}
@outputs^std::domain_error~
\end{example}

\noindent
\code{gaol\_ieee1788::textToInterval()} throws neither: it gives the empty set
for every text that is no interval, as IEEE 1788-2015 has it, so that a
malformed text cannot be told from \code{"[empty]"}:

\begin{example}
cout << gaol_ieee1788::textToInterval("[1, 2") << ' '
     << gaol_ieee1788::textToInterval("[empty]");
@outputs^[empty] [empty]~
\end{example}

\noindent
A null pointer is no text, and the \code{textToInterval()} of both namespaces
and \code{exactToInterval()} throw \code{std::invalid\_argument} for it: none
of the three calls below reaches the reader.

\begin{example}
textToInterval(nullptr);
textToInterval(std::getenv("UNSET"));
gaol_ieee1788::exactToInterval(nullptr);
\end{example}

\noindent
Their argument is a \code{const std::string\&}, which the call builds from the
pointer first, and the C++ standard leaves that undefined (libstdc++ throws
\code{std::logic\_error}); the overloads of a C string take the null pointer
before.\newinvfive

\medskip
Two more texts deserve a word. A sign next to an uncertain number,
\code{10-3.56?1} or \code{10+3.56?1}, throws \code{std::invalid\_argument}: the
lexer reads \code{-3.56?1} (and \code{+3.56?1}) as one token, and
\code{10 -3.56?1} throws too. A blank is needed after the sign:
\code{10 - 3.56?1}, or \code{10- 3.56?1}. Reading a long number takes a time
quadratic in its number of digits: \code{"0."} followed by 20\,000 digits
\code{3} takes 0.017 s, 80\,000 digits 0.28 s, 320\,000 digits 4.4 s (about
45 s expected for a million), and the uncertain form 2.3 times as long. A text
coming from elsewhere can thus block a program.

\subsection{Threads and the global settings}
%%..........................................
\cindex{threads!output of intervals}%
\cindex{interval::format@\code{interval::format}!threads}%

The output format of intervals and their number of digits, set by
\code{interval::\allowbreak format()} and \code{interval::\allowbreak precision()}, are variables of
the program, shared by its threads, and not atomic: setting one while another
thread writes an interval, with \code{operator<{}<}, the conversion to a
\code{std::string} or \code{gaol\_ieee1788::intervalToText()} (which reads
the precision only), is a data race, with the undefined behaviour of C++
(ThreadSanitizer reports it). The program that sets them synchronizes with
the threads that write: it sets them before starting its threads, or takes a
lock of its own around the setting and around each output. Without a setting
in between, several threads write intervals at once safely, and so they read
strings (\pxref{sec:reading-intervals}); \code{exact\_string()} reads neither
setting.\newinvfive

\subsection{The flags of interval arithmetic}
%%...........................................

Compiled without the flags of interval arithmetic (\pxref{sec:compiling}), the
inline operations of GAOL may be evaluated by the compiler at compile time,
rounded to nearest, or contracted into fused multiply-adds, and give bounds
that do not enclose the exact results. \file{gaol/gaol\_config.h} stops the
compilation where it can tell:

\begin{onscreen}
% c++ -ffast-math -c program.cpp
gaol/gaol_config.h: error: #error "GAOL cannot be
compiled with -ffast-math (nor -Ofast): the bounds it
computes would not enclose the exact results (remove
the option, or put the flags of gaol.pc and gaol::gaol
after it)"
\end{onscreen}

\noindent
It stops \option{-ffinite-math-only} in the same way, with which the compiler
takes NaN, the bounds of the empty interval, and infinities never to occur
(\pxref{sec:refused}). Taking the flags from \code{gaol::gaol} or
\file{gaol.pc} avoids the question. The link is not seen by the header: a
program linked with \option{-Ofast} gets the flush-to-zero mode, which makes
the bounds of the operations with subnormal operands and results wrong where
GAOL does not clear it
(\pxref{sec:ftz}).\newinvfive

\section{Coming from another library}\label{sec:other-libraries}
%%--------------------------------------
\cindex{libraries!names in other}%
\cindex{IBEX}%
\cindex{Codac}%
\cindex{C-XSC}%
\cindex{PROFIL/BIAS}%
\cindex{Boost.Interval}%

The names of interval libraries differ, and the same symbol can mean another
relation. Tables~\ref{tab:other-names-ibex} and \ref{tab:other-names-cxsc}
give, for what an algorithm uses, the name of GAOL v5, the one of
\code{gaol\_ieee1788}, and those of IBEX and Codac, which share them, of
C-XSC and of Boost.Interval. A dash means that the sources read have no
function of the kind, not that the algorithm cannot be written.\newinvfive

{\renewcommand{\arraystretch}{1.1}\footnotesize\setlength{\tabcolsep}{3pt}
\begin{longtable}{L{.17\textwidth}L{.24\textwidth}%
    L{.22\textwidth}L{.29\textwidth}}
    \caption{The names of GAOL v5, of \code{gaol\_ieee1788}, and of IBEX and
      Codac}\label{tab:other-names-ibex}\\
    \hline
    & GAOL v5 & \code{gaol\_ieee1788} & IBEX, Codac\\
    \hline\hline
    \endfirsthead
    \hline
    & GAOL v5 & \code{gaol\_ieee1788} & IBEX, Codac\\
    \hline\hline
    \endhead
    Lower bound & \code{x.left()} & \code{inf(x)} & \code{x.lb()}\\
    Upper bound & \code{x.right()} & \code{sup(x)} & \code{x.ub()}\\
    Midpoint, a double & \code{x.midpoint()} & \code{mid(x)} & \code{x.mid()}\\
    Midpoint, an interval & \code{x.mid()} & --- & ---\\
    Width &
    \code{x.width()}, rounded upward; \code{x.width\_enclosure()} &
    \code{wid(x)} &
    \code{x.diam()}\\
    Radius & \code{x.rad()} & \code{rad(x)} & \code{x.rad()}\\
    Smallest, greatest absolute value &
    \code{x.mig()}, \code{x.mag()} &
    \code{mig(x)}, \code{mag(x)} &
    \code{x.mig()}, \code{x.mag()}\\
    Hull, intersection &
    \code{x | y}, \code{x \& y}; \code{hull(x, y)}, \code{intersect(x, y)} &
    \code{convexHull(x, y)}, \code{intersection(x, y)} &
    \code{x | y}, \code{x \& y}\\
    Is empty & \code{x.is\_empty()} & \code{isEmpty(x)} & \code{x.is\_empty()}\\
    Subset &
    \code{x.set\_leq(y)} &
    \code{subset(x, y)} &
    \code{x.is\_subset(y)}\\
    Interior &
    \code{x.set\_le(y)}; \code{y.set\_strictly\_contains(x)} &
    \code{interior(x, y)} &
    \code{x.is\_interior\_subset(y)}\\
    Contains a number &
    \code{x.set\_contains(d)} &
    \code{isMember(d, x)} &
    \code{x.contains(d)}\\
    Disjoint &
    \code{x.set\_disjoint(y)} &
    \code{disjoint(x, y)} &
    \code{x.is\_disjoint(y)}\\
    Equal & \code{x.set\_eq(y)} & \code{equal(x, y)} & \code{x == y}\\
    Inflate by \code{r} & \code{x.inflate(r)} & --- & \code{x.inflate(r)}\\
    Bisect &
    \code{x.bisect(ratio)} &
    --- &
    \code{x.bisect(ratio)}, ratio 0.5 (IBEX) or 0.49 (Codac)\\
    Two-piece division &
    two \code{div\_rel}, or \code{mulRevToPair(b, c)} &
    \code{mulRevToPair(b, c)} &
    \code{div2} (IBEX only)\\
    \hline
  \end{longtable}}

{\renewcommand{\arraystretch}{1.1}\footnotesize\setlength{\tabcolsep}{3pt}
\begin{longtable}{L{.19\textwidth}L{.25\textwidth}%
    L{.25\textwidth}L{.22\textwidth}}
    \caption{The names of GAOL v5, of C-XSC and of
      Boost.Interval}\label{tab:other-names-cxsc}\\
    \hline
    & GAOL v5 & C-XSC & Boost.Interval\\
    \hline\hline
    \endfirsthead
    \hline
    & GAOL v5 & C-XSC & Boost.Interval\\
    \hline\hline
    \endhead
    Lower bound & \code{x.left()} & \code{Inf(x)} & \code{lower(x)}\\
    Upper bound & \code{x.right()} & \code{Sup(x)} & \code{upper(x)}\\
    Midpoint, a double &
    \code{x.midpoint()} &
    \code{mid(x)} &
    \code{median(x)}\\
    Width, rounded upward &
    \code{x.width()} &
    \code{diam(x)} &
    \code{width(x)}\\
    Radius & \code{x.rad()} & --- & ---\\
    Smallest, greatest absolute value &
    \code{x.mig()}, \code{x.mag()} &
    \code{AbsMin(x)}, \code{AbsMax(x)} &
    \code{norm(x)}, the greatest only\\
    Hull, intersection &
    \code{x | y}, \code{x \& y}; \code{hull(x, y)}, \code{intersect(x, y)} &
    \code{x | y}, \code{x \& y} &
    \code{hull(x, y)}, \code{intersect(x, y)}\\
    Is empty & \code{x.is\_empty()} & \code{IsEmpty(x)} & \code{empty(x)}\\
    Subset & \code{x.set\_leq(y)} & \code{x <= y} & \code{subset(x, y)}\\
    Interior & \code{x.set\_le(y)} & \code{x < y}, \code{in(x, y)} & ---\\
    Contains a number &
    \code{x.set\_contains(d)} &
    \code{d <= x}, \code{in(d, x)} &
    \code{in(d, x)}\\
    Disjoint &
    \code{x.set\_disjoint(y)} &
    \code{Disjoint(x, y)} &
    \code{!overlap(x, y)}\\
    Equal & \code{x.set\_eq(y)} & \code{x == y} & \code{equal(x, y)}\\
    Inflate by \code{r} &
    \code{x.inflate(r)} &
    \code{Blow(x, eps)}, relative to the width &
    \code{widen(x, r)}\\
    Bisect &
    \code{x.bisect(ratio)} &
    --- &
    \code{bisect(x)}, at the median only\\
    \hline
  \end{longtable}}

\noindent
\code{overlap(x, y)} of Boost.Interval is a boolean, true when $x$ and $y$
have a point in common: \code{!x.set\_disjoint(y)} in GAOL v5. The
\code{overlap} of GAOL v5 (\code{x.overlap(y)},
\code{gaol\_ieee1788::overlap(x, y)}) is the function of IEEE 1788-2015 and
returns an \code{overlapping\_state}, which does not convert to \code{bool}: a
Boost program that tests \code{overlap(x, y)} does not compile with it.

\medskip
What the symbols \code{<}, \code{<=} and \code{==} mean is where code ported
from another library compiles and changes meaning
(Table~\ref{tab:other-symbols}).

{\renewcommand{\arraystretch}{1.1}\footnotesize\setlength{\tabcolsep}{3pt}
\begin{longtable}{L{.17\textwidth}L{.27\textwidth}%
    L{.27\textwidth}L{.19\textwidth}}
    \caption{The meaning of the symbols \code{<}, \code{<=} and \code{==} in
      each library}\label{tab:other-symbols}\\
    \hline
    & \code{x < y} & \code{x <= y} & \code{x == y}\\
    \hline\hline
    \endfirsthead
    \hline
    & \code{x < y} & \code{x <= y} & \code{x == y}\\
    \hline\hline
    \endhead
    GAOL v5 &
    every point of \code{x} is below every point of \code{y}
    (\ieee{strictPrecedes}) &
    every point of \code{x} is at most every point of \code{y}
    (\ieee{precedes}) &
    not defined: \code{x.set\_eq(y)}\\
    IBEX &
    not defined &
    not defined &
    the same bounds, the empty set equal to itself\\
    Codac &
    a \code{BoolInterval}: true, false, unknown, or empty if an operand is &
    not defined &
    the same bounds, as IBEX\\
    C-XSC &
    \code{x} in the interior of \code{y} &
    \code{x} is a subset of \code{y}; \code{r <= x}: the number \code{r} is in
    \code{x} &
    the same bounds\\
    PROFIL/BIAS &
    \code{x} in the interior of \code{y} &
    \code{x} is a subset of \code{y}; \code{r <= x}: \code{r} is in \code{x} &
    the same bounds\\
    Boost.Interval &
    true if every point of \code{x} is below every point of \code{y}, false if
    none is, else it throws \code{comparison\_error} &
    true if every point of \code{x} is at most every point of \code{y}, false
    if none is, else it throws &
    true if \code{x} and \code{y} are the same point, false if they are
    disjoint, else it throws\\
    \hline
  \end{longtable}}

\noindent
So \code{0.0 <= f} is ``0 is in \code{f}'' in C-XSC and in PROFIL/BIAS, and in
GAOL v5 the certainly relation ``0 is at most every point of \code{f}'': a
program ported from them compiles and no longer tests what it did
(\code{f.set\_contains(0.0)} does). A Boost.Interval program that catches
\code{comparison\_error} to detect an uncertain comparison never sees it with
GAOL v5, whose \code{x < y} is false for overlapping intervals, and the empty
set, which Boost refuses, passes every certainly relation
(\pxref{sec:lexicographic-less}). The names are those of IBEX~2.9.1
(\file{ibex\_Interval.h}), Codac~2 (\file{codac2\_Interval.h}), C-XSC~2.5.4
(\file{interval.hpp} and \file{interval.inl}), PROFIL/BIAS~2.0.8
(\file{Interval.h}) and Boost~1.86 (\file{boost/numeric/interval/}), read in
their sources.

\section{Discussion}
%%------------------

This chapter has shown what a program using GAOL looks like, and what it owes
to the library: one header, the flags of interval arithmetic on every file
that includes it, which CMake, pkg-config and meson pass on, and a call to
\code{gaol::cleanup()} after the last use. The flags are the price of inline
operations: a compiler left free to round to nearest, to fold a constant or
to fuse a product and a sum would break the enclosure, and
\file{gaol/gaol\_config.h} refuses the options it can see, not the absence of
those it cannot. The two namespaces are the answer to two audiences, the
programs written for GAOL~4, IBEX and Codac, and those written to IEEE
1788-2015: the same type and the same functions, under two sets of names,
with the power as their one real difference.

\medskip
The common errors reviewed here all come from the same source: a double is
not an interval, and an interval is not a number. A decimal constant, a
function of a number, a midpoint and a radius, the sign of a zero bound, the
empty set that passes every test, the relations that order nothing, are
each a place where the habits of floating-point programming give a wrong
program that compiles. The tables of names of the other libraries serve the
same purpose for a program ported to GAOL.

\medskip
The next chapter looks at what this one took for granted: the rounding
direction that each operation sets, what it does to the rest of the program,
and the initialization and the cleanup of GAOL.
